\section{面积、体积、概率与累积量}
\label{sec:02-Calculus/04-Integration/05}

\subsection{不止是``曲线下的面积''}

学完定积分的时候，你已经掌握了计算单一曲线下方面积的工具。\(\int_a^b f(x)dx\) 给你 \(y=f(x)\) 从 \(a\) 到 \(b\) 下方那一整块的面积。但真实问题很少这么规矩。

比如两条曲线 \(y=x\) 和 \(y=x^2\)，它们从 \(x=0\) 到 \(x=1\) 夹出了一块区域。这块区域的面积是多少？单独用 \(\int_0^1 x\,dx\) 得到三角形面积 \(1/2\)，单独用 \(\int_0^1 x^2\,dx\) 得到抛物线下方面积 \(1/3\)。但我们要的不是某个曲线下方的面积，而是两条曲线之间的面积。

直觉上，这应该就是大的减去小的：\(1/2-1/3=1/6\)。这个直觉是对的，但只因为 \(x\ge x^2\) 在 \([0,1]\) 上始终成立。如果两条曲线在这个区间内交叉了呢？我们就不能简单做减法了——必须先找出交点，判断各段上谁在上谁在下，分段处理，或直接积分绝对差。

举个交叉的例子立刻能感受到区别。取 \(y=x^3\) 和 \(y=x\) 在 \([-1,1]\) 上。两条曲线的交点在 \(x=-1,0,1\)。\((x^3)-x=x(x^2-1)=x(x-1)(x+1)\)。在 \([-1,0]\) 上，取 \(x=-0.5\)，\((-0.5)^3=-0.125\)，而 \(-0.5\) 更小，所以 \(x^3 > x\)，上方是 \(x^3\)。在 \([0,1]\) 上，\(x=0.5\) 代入：\(0.125 < 0.5\)，所以 \(x > x^3\)。直接做 \(\int_{-1}^1(x^3-x)dx\) 会得到零——因为奇函数的对称性使两段面积抵消。正确的几何面积是：
\[
\int_{-1}^0(x^3-x)\,dx + \int_0^1(x-x^3)\,dx
=\left[\frac{x^4}{4}-\frac{x^2}{2}\right]_{-1}^0+\left[\frac{x^2}{2}-\frac{x^4}{4}\right]_0^1
=\frac14+\frac14=\frac12.
\]

这个具体的面积问题暴露了一个通用的建模思路：积分应用千变万化，但底层动作始终是同一个——把对象切成许多小块，每块的贡献近似写为``局部某种量 \(\times\) 小尺度''，然后把所有小块加总取极限。

\subsection{切片是统一动作，公式只是切法的名字}

无论你在算什么——面积、体积、质量、功、概率——第一步始终是问三个问题：
\begin{enumerate}
\tightlist
\item 切什么？沿着哪个变量把对象切成小片？
\item 每片的近似贡献是什么？用该点的相关信息乘以那个小尺度。
\item 边界从哪里到哪？交点、题目给出的范围、物理约束决定了积分的起止点。
\end{enumerate}

这一套动作下来，积分公式不是要背的，而是自然出现的。``上减下''不是死记的规则——竖直切条时，高就是上边界的 \(y\) 减下边界的 \(y\)。``外减内''也一样——旋转一个区域时，如果区域有内外半径，环的面积就是外圆面积减内圆面积。

\subsection{两条曲线之间的面积：高是上减下}

正式写下来：若在 \([a,b]\) 上恒有 \(f(x)\ge g(x)\)，则面积
\[
A=\int_a^b\bigl[f(x)-g(x)\bigr]\,dx.
\]

每一条竖直小条的宽为 \(dx\)，高为 \(f(x)-g(x)\)，面积微元就是高度乘宽度。所有小条从 \(x=a\) 排到 \(x=b\)，积分把它们累加。

如果上下关系在某段区间发生了变化，简单做 \(\int(f-g)\) 会得到带符号的差，正面积和负面积相互抵消。定积分的定义本身就带符号，所以几何面积要求你在变号处分段取绝对值。

以 \(y=x\) 和 \(y=x^2\) 为例，在 \([0,1]\) 上交点只有 \(x=0\) 和 \(x=1\)。任取中间一点 \(x=0.5\)，\(0.5>0.25\)，所以 \(x\ge x^2\) 在整个区间上成立，不需要分段：
\[
\int_0^1(x-x^2)\,dx
=\left[\frac{x^2}{2}-\frac{x^3}{3}\right]_0^1
=\frac12-\frac13=\frac16.
\]

如果曲线是 \(y=\sin x\) 和 \(y=0\) 在 \([0,2\pi]\) 上，\(\sin x\) 在 \((\pi,2\pi)\) 为负，直接积分会得到零——上半部分的面积被下半部分抵消了。正确的几何面积是 \(\int_0^\pi\sin x\,dx+\int_\pi^{2\pi}(-\sin x)\,dx\)。

\subsection{旋转体的体积：不同切法对应不同公式}

同样是切片，切的方向不同会引出不同的体积公式。

把 \(y=f(x)\ge0\) 在 \([a,b]\) 上的区域绕 \(x\) 轴旋转。沿 \(x\) 方向竖直切：每条竖直小条旋转后变成一个薄圆盘，半径是 \(f(x)\)，厚度 \(dx\)。圆盘的体积微元是底面积 \(\pi[f(x)]^2\) 乘厚度 \(dx\)。因此
\[
V=\pi\int_a^b[f(x)]^2\,dx.
\]

如果区域有内外边界——比如两条曲线 \(y=f(x)\) 和 \(y=g(x)\) 之间的区域绕 \(x\) 轴旋转——每个小片是一个有孔的圆环（垫圈），截面积为 \(\pi([f(x)]^2-[g(x)]^2)\)。

用具体例子试：\(y=\sqrt{x}\)，\(0\le x\le1\)，绕 \(x\) 轴旋转。
\[
V=\pi\int_0^1(\sqrt{x})^2\,dx=\pi\int_0^1 x\,dx=\pi\cdot\frac12=\frac\pi2.
\]

现在同一个区域，改绕 \(y\) 轴旋转。这次沿 \(x\) 方向的竖条不再保持圆柱形状——每个竖直小条旋转后变成的是一个薄圆管，不是实心盘。所以我们改换切法：沿着平行于旋转轴的方向，把区域看成一层层同心圆柱壳的叠加。每个圆柱壳对应一个 \(x\) 值，壳的半径为 \(x\)，周长 \(2\pi x\)，高度为 \(\sqrt{x}\)，壳壁厚度 \(dx\)。剥开这个圆柱壳摊平，近似是一块长 \(2\pi x\)、宽 \(\sqrt{x}\)、厚 \(dx\) 的薄片，体积微元是 \(2\pi x\cdot\sqrt{x}\cdot dx\)：
\[
V=2\pi\int_0^1 x\sqrt{x}\,dx=2\pi\int_0^1 x^{3/2}\,dx=2\pi\cdot\frac25=\frac{4\pi}{5}.
\]

用同一个区域、绕不同轴旋转得到了不同的体积（\(\pi/2\) 对 \(4\pi/5\)），这很合理——旋转轴离区域越远，扫出的体积越大，壳层法里乘的那个 \(2\pi x\) 就是半径的累积效应。

圆盘法和壳层法不是两个不同的理论——它们是同一切片原则在不同切法上的体现。选择哪种，看区域描述在哪个方向上更自然、积分更好算。大致规律：垂直于旋转轴的切片给出圆盘/垫圈；平行于旋转轴的切片给出圆柱壳。如果区域到旋转轴的距离用 \(y\) 描述更简单（如 \(y=f(x)\) 绕 \(x\) 轴），用圆盘；如果用 \(x\) 描述更简单（如 \(y=f(x)\) 绕 \(y\) 轴），用壳层。

\subsection{弧长公式也是同一切片逻辑}

曲线从 \(x=a\) 到 \(x=b\)，把曲线切成许多小段，每一小段近似为直角三角形的斜边：水平步长 \(dx\)，垂直步长 \(dy=f'(x)dx\)，斜边长度 \(\sqrt{(dx)^2+(dy)^2}=\sqrt{1+[f'(x)]^2}\,dx\)。累加所有小段：
\[
L=\int_a^b\sqrt{1+[f'(x)]^2}\,dx.
\]

推导弧长公式不需要新原理——它就是``局部近似长度 \(\times\) 小步长''的同一个框架，只不过这次每小片贡献的不是面积，而是一段长度。

\subsection{密度积分：非均匀材料中的总量恢复}

一根细棒从 \(x=a\) 到 \(x=b\)，线密度随位置变化为 \(\rho(x)\)（单位：千克/米）。密度不均匀意味着不能用``总质量 = 平均密度 \(\times\) 长度''一步得出——因为平均密度本身就是不知道的。

沿长度方向把小棒切成许多长度为 \(dx\) 的小段。在足够短的一段上，密度近似均匀，质量约为 \(\rho(x)\,dx\)。把各段加起来：
\[
M=\int_a^b\rho(x)\,dx.
\]

同样的逻辑可以求质心——质心是位置按质量加权的平均。位置 \(x\) 处的小段质量为 \(\rho(x)dx\)，它对该段位置的``贡献权重''是 \(x\cdot\rho(x)dx\)。所有段的一阶矩之和除以总质量：
\[
\bar x=\frac1M\int_a^b x\,\rho(x)\,dx.
\]

分母 \(M\) 的作用是归一化——如果直接用一阶矩而不除总质量，你会得到一个依赖于总质量的量而非纯位置。当密度均匀时 \(\rho\) 为常数，\(\bar x=(a+b)/2\)，恢复几何中点。密度偏向右边时，质心也右移。

\subsection{变力做功：不是力乘距离，而是力乘每一小段距离再累加}

物理课上学过恒力做功：\(W=F\cdot d\)。但如果力随位置变化呢？拉一根弹簧，拉力 \(F(x)=kx\) 随伸长量线性增加。不能直接用``某个力乘总距离''。

把位移从 \(x=0\) 到 \(x=X\) 切成许多小区间。在 \([x,x+dx]\) 这一小段上，力近似不变，做的功约为 \(F(x)\,dx\)。总功是所有小段的累加：
\[
W=\int_0^X kx\,dx=\frac12 kX^2.
\]

弹簧的弹性势能公式 \(\frac12 kX^2\) 和匀加速直线运动的位移公式 \(\frac12 a t^2\) 结构一致——都来自线性增长的累积。``切段、近似、累加、取极限''这个流程，一旦习惯，你会发现它在物理、工程、统计里无处不在。

\subsection{概率密度：注意单位}

连续随机变量的密度函数 \(p(x)\) 进入积分框架的方式和物理密度高度相似，但有一个极易误解的单位差异。

\(p(x)\) 满足：
\[
p(x)\ge0,\qquad \int_{-\infty}^{\infty}p(x)\,dx=1.
\]

区间概率定义为
\[
P(a\le X\le b)=\int_a^b p(x)\,dx.
\]

这里的关键点是：\(p(x)\) 的单位是``每单位 \(x\) 的概率''。物理密度是``每单位长度的质量''，概率密度是``每单位 \(x\) 的概率''——结构完全一致。正因为是密度，\(p(x)\) 的数值可以大于 1。例如在 \([0,0.1]\) 上密度恒为 10 的均匀分布完全合法——面积仍是 \(0.1\times10=1\)。单个点的概率是零，因为零宽的区间没有面积。但这不意味着该点``不可能发生''，只意味着你无法在连续尺度上给单个点分配非零概率。

累积分布函数
\[
F(x)=P(X\le x)=\int_{-\infty}^x p(t)\,dt
\]
从左侧把所有概率密度一路扫过来。在密度连续的点上，微积分基本定理给出 \(F'(x)=p(x)\)：累积量的瞬时变化率就是密度本身。

\subsection{平均值是归一化后的累积}

函数在区间上的平均值：
\[
f_{\mathrm{avg}}=\frac1{b-a}\int_a^b f(x)\,dx.
\]

这个式子的结构跟质心一致：\(\int_a^b f(x)\,dx\) 是 \(f\) 在区间上的累积，除以区间长度 \(b-a\) 做归一化。若 \(f\) 恒为常数，累积量就是常数乘长度，除以长度还原常数本身。

积分中值定理保证：若 \(f\) 在 \([a,b]\) 上连续，则存在某个 \(c\in[a,b]\) 使得 \(f(c)=f_{\mathrm{avg}}\)。推导思路很直接：平均值介于 \(f\) 的最小值与最大值之间，而连续函数的值域是一个区间——介值定理保证最小值和最大值之间的每个值都会被取到，所以平均值必定是某点的函数值。

但平均值不是简单把端点值取平均——除非 \(f\) 刚好是线性函数。``中间值''指的是函数值会经过平均高度，不是自变量取在区间正中。

\subsection{建模才是功夫所在}

上面这些应用套到一个共同的结论：积分本身是最后一步，把物理或几何情境转换成积分表达式才是真正需要判断力的地方。你一旦确定了切的方向、每片的近似贡献、积分的起点和终点，剩下的定积分计算——大部分情形下——可以用基本定理或数值方法完成。积分的应用不是在考你积分技术，而是在考你能否把一种连续的累加过程精确地写成数学语言。

因此，在动手算之前先画图并标注变量，确保你清楚：切片沿哪个方向走，边界在哪里，哪条曲线在上。这三个问题答对了，式子就不会错。

\subsection{延伸与来源}

\begin{itemize}
\tightlist
\item
  MIT OpenCourseWare 18.01SC：The Definite Integral and its Applications
\item
  MIT OpenCourseWare 18.01SC：Techniques of Integration
\end{itemize}
